\subsection{从事件到条件概率}

基本事件是一次试验不可再细分的结果；不同基本事件互斥，所有基本事件的并集是样本空间。
概率满足 $P(A)\ge0$、$P(\Omega)=1$。若 $A,B$ 互斥，$P(A\cup B)=P(A)+P(B)$；
一般情形则须扣除重复计算的交集：$P(A\cup B)=P(A)+P(B)-P(A\cap B)$。
有限个基本事件等可能时，概率等于有利结果数除以总结果数，这叫古典概率模型。
例如公平骰子掷出偶数的概率是 $3/6=1/2$。连续区域内均匀取点时，概率等于长度、面积或体积之比，
这是几何概率模型；必须先说明对哪一种几何测度均匀。
连续样本空间的事件须可测。例如在 $[0,1]$ 上均匀取点，“取到有理数”是可测事件，概率为零，
因为有理数可数且每个点的概率为零；概率为零不等于集合为空。某些不可测集合则不能在该测度下直接赋予概率。
多次试验的联合结果是一个数组 $(\omega_1,\ldots,\omega_n)$，联合样本空间是各空间的直积；
只有另加独立性假设，联合概率才是各分量概率的乘积。

已经知道事件 $A$ 发生，相当于把样本空间缩小到 $A$，并重新归一化。只要 $P(A)>0$，
\begin{equation}
 P(B\mid A)=\frac{P(A\cap B)}{P(A)},\qquad
 P(A\cap B)=P(A)P(B\mid A).
\end{equation}
已知骰子点数为偶数，再问点数是否为6，概率变为 $(1/6)/(1/2)=1/3$，不是原先的 $1/6$。
独立的含义是知道 $A$ 不改变 $B$ 的概率，即 $P(B\mid A)=P(B)$，从而得到乘法公式。
“互斥”与“独立”不同：两个正概率的互斥事件不可能独立，因为一个发生时另一个的条件概率是零。
独立性对双方是对称的，也传给补事件。例如
$P(A\cap B^c)=P(A)-P(A\cap B)=P(A)[1-P(B)]$，
同理 $P(A^c\cap B)=P(A^c)P(B)$、$P(A^c\cap B^c)=P(A^c)P(B^c)$。
随机变量是样本空间上的实值函数 $X(\omega)$；同一个数值可对应多个基本事件，其概率须将这些事件合并计算。

对于连续变量，$f_X(x)\dd x$ 是落在微小区间内的概率，而 $f_X(x)$ 本身不是概率；它的单位是 $X$ 的单位的倒数。
\[P(a<X<b)=\int_a^bf_X(x)\dd x,\qquad
f_X(x)\ge0,\qquad\int f_X(x)\dd x=1.\]
密度可以大于1，但任一区间的概率不能大于1。两个变量的联合密度 $f_{XY}$ 与边缘密度的关系为
\[f_X(x)=\int f_{XY}(x,y)\dd y,\qquad
f_Y(y)=\int f_{XY}(x,y)\dd x.\]
独立时联合密度才可分解为 $f_X(x)f_Y(y)$。离散情形对应 $p_i=\sum_jp_{ij}$、$q_j=\sum_ip_{ij}$。

\subsection{变量变换为什么要累加不同原像}

设无量纲速度分量 $u$ 服从标准高斯密度
$f_u(u)=(2\pi)^{-1/2}\exp(-u^2/2)$，令无量纲动能 $e=u^2/2$。
任意 $e>0$ 都有两个原像 $u_+=\sqrt{2e}$、$u_-=-\sqrt{2e}$。
两个速度区间都贡献同一能量区间的概率，因此
\begin{align*}
 f_e(e)&=f_u(u_+)\left|\frac{\dd u_+}{\dd e}\right|
       +f_u(u_-)\left|\frac{\dd u_-}{\dd e}\right|\\
 &=2\frac{e^{-e}}{\sqrt{2\pi}}\frac1{\sqrt{2e}}
 =\frac{e^{-e}}{\sqrt{\pi e}},\qquad e>0.
\end{align*}
代换 $e=s^2$ 检验归一化：
\[\int_0^\infty f_e(e)\dd e
=\frac2{\sqrt\pi}\int_0^\infty e^{-s^2}\dd s=1.\]
密度在零点附近发散但可积，并不表示零能量具有有限点概率。若只取正根，积分会错成 $1/2$。
一般变换公式在各单调分支上使用，导数为零的端点通过极限处理。

\subsection{样本平均与宏观稳定性}

设独立同分布变量的均值为 $\mu$、有限方差为 $\sigma^2$，样本平均为 $\bar X=N^{-1}\sum_iX_i$。
线性性给出 $\langle\bar X\rangle=\mu$；独立性使交叉协方差为零，于是
\[\operatorname{Var}(\bar X)=\frac1{N^2}\sum_i\operatorname{Var}(X_i)
=\frac{\sigma^2}{N}.\]
对于任何固定容许误差 $\delta>0$，切比雪夫不等式给出
\[P(|\bar X-\mu|\ge\delta)\le\frac{\sigma^2}{N\delta^2}\longrightarrow0.\]
这证明了有限方差情形的弱大数定律：样本平均依概率趋向总体平均，而不是每一次有限试验必然完全相等。
期望或方差不存在时不能直接套用本推导。例如柯西密度 $1/[\pi(1+x^2)]$ 的通常期望积分不收敛，
不能因图形关于零对称就把其期望当作零。若变量相关，方差中还须保留 $2\sum_{i<j}\operatorname{Cov}(X_i,X_j)$。
